\scanpage{519}
\mspage{430}
% B01
à prouver que pour tout $\lambda:s\relbar t$, la flèche
% B02
\[\underset i{\unsure{\underline{\lim}|\varinjlim}}\,\mathrm{Hom}_{\lambda}(x,y_i)\longrightarrow\mathrm{Hom}_{\lambda}(x,y)\]
% B03
est bijective. Pour \unsure{$a$|$x$} \inserted{et $\lambda$} \deleted{\unsure{$a$|$x$}} fixés, et des syst. inductifs \unsure{variables|arbitraires} dans $E_t$, grd devant $\pi$, cela signifie que le foncteur
% B04
\[H_\lambda(x)=(y\longmapsto\mathrm{Hom}_\lambda(x,y)):\quad E_t\unsure{\relbar|\longrightarrow}\mathrm{Ens}\]
% B05
est $\pi$-accessible, et ceci quel que soit
% B06
\[x\in\mathrm{Ob}\,E_{s\pi}.\]

% B07
Or c'est le cas si $x\in\mathrm{Ob}\,E_{s\pi_0}$, en vertu de l'hyp. $3^\circ$) (\deleted{\unsure{car que}} alors $H_\lambda(x)$ est $\pi_1$-accessible, et a fortiori $\pi$-accessible). Il faut \deleted{\unsure{prouver}} passer de là au cas où $x\in\mathrm{Ob}\,E_{s\pi}$. Or comme $\pi_0$ est utile pour $E_s$ par $1^\circ$, on a
% B08
\[x=\varinjlim_{j\in J}x_j\qquad\left|\begin{aligned}&x_j\in\mathrm{Ob}\,E_{s\pi_0}\\&J\ \text{grd devant }\pi_0,\end{aligned}\right.\]
% B09
et comme \inserted{par $2^\circ$)} \inserted{on a} $\pi^{\pi_0}=\pi\geq\pi_1\geq\deleted{\unsure{\mathrm{card}}}\mathrm{Sup}(\pi_0,\mathrm{card\,red}\ E_{s\pi_0})$, on en conclut qu'on peut prendre $J$ de sorte que
% B10
\[\mathrm{card}\,J\leq\pi.\]

% B11
D'autre part, en vertu de $\pi\geq\pi_1\geq\pi_\lambda$, le foncteur
% B12
\[\left|\begin{aligned}x_0&\longmapsto H_\lambda(x_0)\\ E_s&\longrightarrow(\underline{\mathrm{Hom}}(E_t,\mathrm{Ens}))^\circ\end{aligned}\right.\]
% B13
est $\pi_0$-accessible, donc on a
% B14
\[H_\lambda(x_0)\simeq\varprojlim_{j\in J}H_\lambda(x_j).\]

% archive sheet number: 518
